\chapter{Floating-point reproducibility}\label{ch:repro}

The port requires that the GPU build of WRF give the same bits as a CPU build of the same source. This chapter
explains why that is the right requirement for a fire model, what stands in its way, and how the port meets it and
checks it. The general background is in \citet{goldberg1991} and \citet{muller2018}.

\section{Why bits matter}

Weather is chaotic: small differences in the state grow until they are as large as the natural variability. A
coupled fire model adds a threshold on top of this:
\begin{itemize}
  \item a fire cell is either burning or not;
  \item the front is where a continuous field (the level-set function) changes sign.
\end{itemize}
A difference in the last bit of a wind component can move the front across a cell boundary a little earlier or
later. The fire then spreads from a different set of cells, and from that moment the difference is no longer
small.

Experiment E1 of the port's plan measures this for the case of \cref{ch:case}. It changes the temperature at one
grid point of the initial state by one unit in the last place, and compares the burned areas after 17~hours. The
result will be reported here (card A-09 of the roadmap).

The consequence for a GPU port is simple. A GPU build that rounds differently from the CPU build produces a
different fire. Statistical comparison of such runs cannot tell a correct port that rounds differently from a port
with a small bug. If the two builds agree bit for bit, every difference is a bug, and the first field that differs
points to the routine that causes it. The port therefore defines correctness as bitwise identity with a CPU
reference build (CPU-REF), and treats ``close'' as a failure.

\section{IEEE arithmetic in brief}

WRF computes mostly in IEEE~754 binary32 (single precision, \code{REAL(4)}) \citep{ieee754}:
\begin{itemize}
  \item 1 sign bit, 8 exponent bits and 23 fraction bits;
  \item a significand of $p = 24$ bits, so the unit roundoff is $u = 2^{-24}$.
\end{itemize}
The distance between two consecutive floating-point numbers, the \emph{unit in the last place} $\ulp(x)$, is
$2^{-23}$ for $x \in [1,2)$.

The standard requires the basic operations $+$, $-$, $\times$, $/$ and $\sqrt{\phantom{x}}$ to be \emph{correctly
rounded}. The result is the exact result rounded to a representable number, by default to the nearest one, with ties
to the one whose last bit is even:
\begin{equation}
  \fl(a \circ b) = \operatorname{round}(a \circ b), \qquad \circ \in \{+,-,\times,/\}.
\end{equation}
This makes the basic operations deterministic: the same operands give the same bits on every conforming processor,
CPU or GPU.

It does not make arithmetic associative. In binary32, $2^{24} = 16\,777\,216$ and the next number is $2^{24}+2$, so
\begin{align}
  \fl(\fl(2^{24} + 1) + 1) &= \fl(2^{24} + 1) = 2^{24} \quad\text{(a tie, rounded to even), but}\\
  \fl(2^{24} + \fl(1 + 1)) &= 2^{24} + 2.
\end{align}

\begin{keypoint}
A computation is reproducible when the same operations are applied in the same order to the same operands. Every
source of irreproducibility breaks one of these three.
\end{keypoint}

\section{What makes CPU and GPU results differ}

\Cref{tab:repro-sources} lists what can make a GPU build differ from a CPU build of the same source, what the port
does about each, and the test that checks it. The sections below explain the main ones.

\begin{table}[ht]
\centering
\small
\begin{tabular}{p{0.24\textwidth}p{0.45\textwidth}p{0.21\textwidth}}
\toprule
Source & Remedy in the port & Test \\
\midrule
Fused multiply-add & contraction off in both builds (\code{-Mnofma}, \code{-gpu=nofma}; gfortran
  \code{-ffp-contract=off}) & T-FMA \\
Elementary functions & own portable implementations, \code{module_repro_math} & T-RM-*, T-SYM \\
Order of operations & statements copied verbatim; sums kept sequential in source order; no reassociation by the
  vectorizer (\code{-Mvect=noassoc}) & T-AB, T-TRACE, \code{arith_guard} \\
Subnormal numbers & no flush to zero (\code{-Mnoflushz}, \code{-gpu=noflushz}) & T-SUBNORM \\
Signed zeros, NaN in \code{MAX}, \code{MIN}, \code{SIGN} & checked; replaced if host and device differ &
  T-SIGNZERO, T-MINMAX \\
Integer powers & checked for each exponent used & T-IPOW \\
Uninitialized memory & scratch and work arrays zero-filled in both builds & T-UNINIT \\
Integer overflow (random numbers) & explicit 32-bit arithmetic & T-KISS \\
Data races & every variable listed (\code{default(none)}); every written scalar private & T-AB, call check \\
\bottomrule
\end{tabular}
\caption{Sources of CPU/GPU differences and how the port handles them.}\label{tab:repro-sources}
\end{table}

\subsection{Fused multiply-add}

A fused multiply-add (FMA) computes $a \times b + c$ with a single rounding. It is more accurate than a multiply
followed by an add, and it is the native operation of GPU floating-point units. Compilers for GPUs use it by default
whenever the source has a product followed by a sum. The result then differs from the CPU result, which rounds twice.

An example in binary32: take $a = b = 1 + 2^{-12}$ and $c = -1$. The exact product is $1 + 2^{-11} + 2^{-24}$. This
lies exactly halfway between the representable numbers $1 + 2^{-11}$ and $1 + 2^{-11} + 2^{-23}$, so it rounds to
the even one, $1 + 2^{-11}$. Therefore
\begin{align}
  \text{unfused:}\quad & \fl(\fl(a b) + c) = 2^{-11},\\
  \text{fused:}\quad & \fl(a b + c) = 2^{-11} + 2^{-24}.
\end{align}
The port's test T-FMA computes exactly this on the host and on the device. It reads the operands from a file at run
time, so that the compiler cannot fold the expression, and it adds a million random triples for which the two
results differ. Contraction is switched off in both builds, and T-FMA must pass before any other GPU test is
believed.

\subsection{Elementary functions}

IEEE~754 recommends, but does not require, correctly rounded elementary functions ($\exp$, $\log$, $\sin$, $x^y$,
\dots). Math libraries trade accuracy for speed, and they differ between vendors, between CPU and GPU, and between
versions. A result that differs in the last bit in one call of \code{EXP} is enough to change the fire.

The port does not call any vendor math library. \code{module_repro_math} implements every elementary function WRF
uses with only the basic operations and exact bit manipulation:
\begin{itemize}
  \item \code{REAL(4)} arguments are converted to double precision and evaluated with a fixed double-precision
    algorithm, whose error is below one ulp of double precision. The result is then rounded once to single
    precision.
  \item \code{REAL(8)} arguments use the algorithms of fdlibm~5.3 \citep{fdlibm}, translated to Fortran.
  \item NaN results are replaced by one canonical quiet NaN, because hardware NaN payloads differ between CPU and
    GPU.
\end{itemize}

Each function is \code{ELEMENTAL} and is compiled for both the host and the device. Since it consists only of basic
operations in a fixed order, it gives the same bits on both. The same module, compiled without \code{REPRO_MATH},
returns the Fortran intrinsics, which gives the upstream behaviour.

\begin{portnote}
\begin{itemize}
  \item \textbf{Correct rounding.} The single-precision functions were tested exhaustively, on all $2^{32}$ inputs
    of each of $\exp$, $\log$, $\log_{10}$, $\sin$, $\cos$, $\tan$, $\arcsin$, $\arccos$, $\arctan$, $\sinh$, $\cosh$
    and $\tanh$. Every result is correctly rounded. $x^y$ is correctly rounded on 7.3 million samples, including
    every exponent WRF uses.

    Correct rounding is a bonus, not the requirement: double rounding (to double, then to single) could in principle
    miss it for an argument extremely close to a rounding boundary, and the result would still be the same on every
    processor.
  \item \textbf{Substitution.} A tool rewrote 646 intrinsic calls and 265 powers in 38 source files into calls of
    the portable functions.
  \item \textbf{Symbol audit.} T-SYM checks that no math-library symbol is left in the compiled objects.
\end{itemize}
\end{portnote}

\subsection{Inside a portable function: the sine}

How can a function be portable and still accurate? \Cref{alg:rp-sin} shows the double-precision sine on which
\code{rp_sin} is built (\src{frame/module_repro_math.F}{rm_sin}, \code{rem_pio2}, \code{k_sin}), translated from
fdlibm \citep{fdlibm}. It has two parts.

\paragraph{Argument reduction.} The argument is written as $x = n\,\pi/2 + y$ with $|y| \le \pi/4$, and
$\sin x$ is then $\pm\sin y$ or $\pm\cos y$, depending on $n \bmod 4$. The difficulty is that $\pi/2$ is not a
floating-point number. Subtracting $n$ times a rounded $\pi/2$ would leave the rounding error of $\pi/2$ multiplied
by $n$ in $y$, and for large $n$ that error would be larger than $y$'s last bit. The algorithm therefore splits
$\pi/2$ into parts (the method of Cody and Waite):
\begin{equation}
  \frac{\pi}{2} = p_1 + p_{1t}, \qquad p_{1t} = p_2 + p_{2t}, \qquad p_{2t} = p_3 + p_{3t},
\end{equation}
where $p_1$, $p_2$ and $p_3$ have only 33 significant bits. For $|n| < 2^{20}$ the product $n\,p_1$ then has at most
53 bits and is computed exactly, so $x - n\,p_1$ is exact as well. Only the small tail $n\,p_{1t}$ is rounded. When
the result has lost many leading bits by cancellation (when $x$ is close to a multiple of $\pi/2$), the next part is
subtracted too, giving about 118 and then 151 correct bits. The reduced argument is returned as a pair $y_0 + y_1$,
its head and its rounding error.

\paragraph{Polynomial.} On $|y| \le \pi/4$, $\sin y$ is a polynomial of degree 13 in $y$, with coefficients
chosen to minimize the maximum error (a minimax polynomial). It is evaluated in Horner form, with the tail $y_1$ of
the reduced argument entering as a first-order correction.

\begin{algorithm}[ht]
\caption{The portable sine (\code{rm_sin}, \code{rem_pio2}, \code{k_sin}; double precision). For \code{REAL(4)},
\code{rp_sin} converts $x$ to double, calls this, and rounds once to single.}
\label{alg:rp-sin}
\begin{algorithmic}[1]
\Require $x$ in double precision
\If{$|x| \le \pi/4$} \Return $k_{\sin}(x, 0)$ \EndIf
\If{$x$ is $\pm\infty$ or NaN} \Return the canonical NaN \EndIf
\If{$|x| < 3\pi/4$} \Comment{$n = \pm1$}
  \State $z \gets |x| - p_1$; \quad $y_0 \gets z - p_{1t}$; \quad $y_1 \gets (z - y_0) - p_{1t}$; \quad $n \gets 1$
\ElsIf{$|x| \le 2^{19}\pi/2$} \Comment{medium arguments}
  \State $n \gets \operatorname{int}(|x|\cdot(2/\pi) + 0.5)$; \quad $r \gets |x| - n\,p_1$ (exact); \quad
    $w \gets n\,p_{1t}$; \quad $y_0 \gets r - w$
  \If{$y_0$ lost more than 16 bits to cancellation}
    \State $t \gets r$; $w \gets n\,p_2$; $r \gets t - w$; $w \gets n\,p_{2t} - ((t - r) - w)$; $y_0 \gets r - w$
    \If{it lost more than 49 bits} \State the same with $p_3$, $p_{3t}$ \EndIf
  \EndIf
  \State $y_1 \gets (r - y_0) - w$
\Else
  \State Payne--Hanek reduction with the bits of $2/\pi$ (\code{k_rem_pio2})
\EndIf
\State restore the sign of $x$ in $y_0$, $y_1$ and $n$
\State \Return $k_{\sin}(y_0, y_1)$, $k_{\cos}(y_0, y_1)$, $-k_{\sin}(y_0,y_1)$ or $-k_{\cos}(y_0,y_1)$ for
  $n \bmod 4 = 0, 1, 2, 3$
\end{algorithmic}
\end{algorithm}

Every step is a basic operation or an exact manipulation of bits, so the result is the same on every IEEE processor,
provided no step is contracted into an FMA. The reduction is a case in point: $t - n\,p_2$ evaluated with an FMA
would round differently from the two roundings the algorithm was designed for. Correct rounding of the final
single-precision result follows from the double-precision accuracy of the steps, and the exhaustive test confirms it
for every one of the $2^{32}$ inputs.

\subsection{Order of operations}

Reassociation changes results, as shown above. The rule of the port is therefore that a GPU kernel performs the
same statements as the CPU loop it replaces, copied verbatim; only the array indices of a slab that becomes a column
may change. Three places need care.
\begin{itemize}
  \item \textbf{Sums.} A sum over species or over levels stays a sequential loop inside one GPU thread, in the
    source order. It is never turned into a parallel reduction, which would add in a different order.
  \item \textbf{Vectorizers.} The CPU compiler may also reorder a reduction to vectorize it. CPU-REF forbids this
    with \code{-Mvect=noassoc}.
  \item \textbf{Integer reductions.} Reductions of integers (counts, flags, checksums) are exact in any order and
    may run in parallel.
\end{itemize}
A static checker, \code{arith_guard}, compares every statement in the GPU code with the CPU lines it replaces.

\subsection{Compiler flags}

The flags of the two builds (\cref{ch:iobuild}) state the rules of this chapter to the compilers:
\begin{center}
\small
\begin{tabular}{lll}
\toprule
Rule & Host (\code{nvfortran}) & Device (\code{-gpu=...}) \\
\midrule
no FMA contraction & \code{-Mnofma} & \code{nofma} \\
no flush of subnormals to zero & \code{-Mnoflushz -Mnodaz} & \code{noflushz} \\
IEEE-conforming division and square root & \code{-Kieee} & (default) \\
no reassociation by the vectorizer & \code{-Mvect=noassoc} & --- \\
own elementary functions & \code{-DREPRO_MATH} & \code{-DREPRO_MATH} \\
\bottomrule
\end{tabular}
\end{center}
Flags are only promises of the compiler. Each is checked by a test that would fail if the promise were broken: T-FMA
for contraction, T-SUBNORM for subnormals and correctly rounded division and square root, T-SYM for leftover
math-library calls in the objects.

\subsection{Rewritten powers}

Some rewrites are exact and are used freely:
\begin{itemize}
  \item \code{x**2.0} becomes \code{x*x}: a correctly rounded square is the correctly rounded product.
  \item \code{x**0.5} becomes \code{SQRT(x)} for the non-negative arguments where WRF uses it.
\end{itemize}

Integer powers are different. A compiler expands \code{x**4} into multiplications, and it may compute
$(x\cdot x)\cdot(x\cdot x)$ or $((x\cdot x)\cdot x)\cdot x$. These can round differently, and host and device
compilers need not choose the same expansion. T-IPOW compares host and device for every integer exponent left in the
code: 2, 3, 4, 8 and 32.

\subsection{Uninitialized memory}

Fortran automatic arrays are not initialized. A CPU build reads whatever the stack contained, and a GPU build reads
something else. A routine that reads an element before writing it therefore behaves differently on the two. The port
replaces the large automatic arrays by work arrays and a scratch pool, which are zero-filled once in both builds.
T-UNINIT checks the rest: a run that initializes every local to a signaling NaN and a run that initializes them to
zero must agree.

\section{What reproducibility costs}

Bitwise reproducibility is not free.
\begin{itemize}
  \item Without FMA, a GPU reaches only half of its peak floating-point rate. Memory-bound code, which is most of
    WRF, loses much less.
  \item The portable elementary functions are slower than the hardware approximations GPUs provide.
  \item Sums inside a column cannot be parallel.
\end{itemize}

The port therefore plans two build modes from one source (ADR-003):
\begin{itemize}
  \item \textbf{REPRO}, bitwise and used for all verification;
  \item \textbf{FAST}, with the vendors' fast arithmetic. It is validated statistically against an ensemble of REPRO
    runs, in the manner of the ensemble consistency tests of \citet{baker2015} and \citet{milroy2018}.
\end{itemize}
The speed difference between the two is the measured cost of reproducibility (\cref{ch:perf}).

\section{How identity is checked}

\paragraph{The bit-hash tracer.} At chosen points of the time step, \code{module_bittrace} reduces each field to
two 64-bit integers computed from the bit patterns of its values over the patch interior:
\begin{equation}
  h_1 = \sum_{g} b_g \bmod 2^{64}, \qquad
  h_2 = \sum_{g} \bigl((b_g \bmod M)\, w_g \bmod M\bigr) \bmod 2^{64}, \qquad M = 2^{31} - 1,
\end{equation}
where $b_g$ is the bit pattern of the value at global point $g$, read as an unsigned 32-bit integer, and
$w_g \in [1, M-1]$ a weight computed from $g$. The first word changes when any bit changes; the second also changes
when values move to other points, which a plain sum would not notice
(\src{frame/module_bittrace.F}{bt_field3}).
The arithmetic is exact integer arithmetic. The hashes therefore do not depend on the order of evaluation, on the
MPI decomposition, or on whether the reduction ran on a CPU or a GPU. Two runs are compared line by line, and the
first differing line names the step, the Runge--Kutta stage, the place in the step and the field.

\paragraph{Tests at three scales.}
\begin{itemize}
  \item \emph{Per routine, device against host} (T-AB). A route switch runs one ported routine on the host and the
    rest on the device. The two traces must agree.
  \item \emph{Per routine, random inputs} (the harness). The routine is called alone with random inputs, in the CPU
    view and in the GPU view. This reaches code paths a short run does not.
  \item \emph{Whole model, GPU against CPU-REF} (T-TRACE). Short windows restarted from the same file, then the full
    17-hour run, compared field by field and fire cell by fire cell.
\end{itemize}

\paragraph{Without a GPU.} Much can be checked on a CPU. gfortran with \code{-fopenacc} compiles the GPU view of the
source and runs its OpenACC regions on the host. The GPU view must then give the same bits as the CPU view. This
proves that the restructured kernels keep the arithmetic: a column that replaced a slab, a flux split into two
kernels, a limiter split into flags and faces.

A mutation test shows that the check has teeth. The correct column port of \code{calc_coef_w} passes, and the same
port with one product reassociated fails. What a CPU cannot check is everything that depends on two memories and
many threads:
\begin{itemize}
  \item data movement between host and device memory, which are the same memory there;
  \item data races, since the host fallback runs each region in a single thread.
\end{itemize}
These are left to the GPU tests: T-AB, the call check, and T-TRACE.

\section{Limits}

Bitwise identity holds between two builds of one source by one compiler family: CPU-REF and GPU-REPRO, both built
with \code{nvfortran}. In REPRO mode only IEEE basic operations and the portable functions remain, so identity across
compilers and vendors should also be within reach (card G-06 of the roadmap). That would mean the same results on
every GPU.

Finally, bitwise identity says only that the port did not change the model. Whether the model forecasts the fire
well is a different question, which this book leaves to the cited literature and to the comparison with
observations.
